\centerline{\bf \Huge ДЗ. Процессы и операции}
\centerline{\bf \Huge }

\begin{flushright}
        \scriptsize{\textbf{Ответы без объяснения решения оцениваются в 0 баллов!}}
\end{flushright}

Каждая задача стоит \textbf{3 балла}, но при этом 1 и 2 задачи \textbf{обязательные}! 

За неправильное выполнение 1 и 2 задачи можно получить \textbf{отрицательное} количество баллов!

\problem{1!} На экране компьютера - число 2031. Каждую секунду число на экране умножают или делят либо на 2, либо на 3. Результат действия возникает на экране вместо записанного числа. Ровно через минуту на экране появилось число. Могло ли это быть число 52? 


\problem{2!} СОЛОИДЖИЛ удалил все натуральные числа, которые делятся на 52. Какое теперь число будет стоять на 525252 месте?

\problem{3} В девятиэтажном доме живет Анджур Аркадьевич и его ученики. Лифт курсирует между вторым и препоследним этажами, останавливаясь на каждом этаже. На каждом этаже, начиная с четвёртого, в лифт заходил один ученик, но никто не выходил. Когда в лифт зашёл 993 ученик, лифт остановился. На каком этаже это произошло?


\problem{4}  КИМЧЕНЫН удалил все натуральные числа, которые делятся на 5 или на 11. Какое теперь число будет стоять на 2024 месте?